\section{Related work}
\paragraph{Curriculum learning.}
Curriculum learning has long studied useful presentation orders \citep{bengio2009}. Transfer-based and convex analyses connect example difficulty to convergence \citep{weinshall2018,weinshall2020}; empirical work shows that benefits depend on the training regime \citep{hacohen2019,wu2021}. Teacher--student analysis separates learning-speed benefits from changes in final performance \citep{saglietti2022}. Task-level curricula also study the order of distinct learning problems \citep{pentina2015}. We do not claim that curricula can help for the first time. Our comparator is the globally best static mixture within a stated dynamics, and our principal obstruction concerns directional initial-error mass.

\paragraph{Data-mixture optimization.}
DoReMi and DoGE learn useful domain weights through proxy optimization \citep{xie2023,fan2024}. Skill-it and Aioli exploit skill structure or update mixture parameters during training \citep{chen2023,chen2025}. RegMix predicts effective static mixtures, while RegMix-D explicitly transfers dynamic schedules from proxy trajectories \citep{liu2025,zhao2026}. Olmix studies mixture design and reuse as the domain set changes \citep{chen2026}. Mixture scaling laws predict losses across source proportions and scales \citep{ye2025,shukor2025}; \citet{dai2026} explain static mixing through capacity competition and noise reduction. Dynamic mixtures, proxy-based scheduling, and theoretical mixture models are therefore established ideas. Our results instead give an exact static benchmark, global schedule-rate comparisons, and a necessary-and-sufficient boundary over all measurable schedules in a deliberately controlled model.

\paragraph{Geometry and control.}
\citet{sweeney2026} uses Lie brackets to predict local transfer order. \citet{mori2025} derive task-selection protocols from exact learning dynamics and optimal control; \citet{gu2025} formulate language-model data selection using Pontryagin conditions. Switched linear systems also have an established control-theoretic treatment \citep{liberzon2003}. A control formulation, a commutator calculation, or a local numerical schedule search is not our contribution. Noncommutativity alone is insufficient for our stronger static-mixture comparison. We use the classical Golden--Thompson inequality \citep{golden1965,thompson1965} and scalar spectral Jensen inequalities explicitly, and do not present their corollaries as new matrix inequalities.

\paragraph{Spectral scaling.}
Solvable regression and random-feature models explain scaling through spectral tails and training dynamics \citep{maloney2022,bordelon2022,bordelon2024}. Feature learning can change that picture \citep{bordelon2025}. Our block result is a comparison statement: under a common spectral envelope and a uniformly adequate static mixture, a shared schedule cannot improve the exponent. It neither replaces those scaling theories nor covers feature learning in a general neural network. Finally, our perturbation and confidence calculations use standard matrix-exponential and Gaussian concentration tools \citep{higham2008,boucheron2013}; the decision being certified is the paid curriculum-versus-full-budget-static comparison.

